% ======================================================================
%  TIER 2 KARTI: API & Authorization
% ======================================================================

\card{API \& Authorization}{eylemin arabirimi ve yetki sınırı}
\cl{Problem} Model bir eylemi tetikler, ama o eylemin teknik arabirimi ve
\emph{kimin neye erişebileceği} koddadır --- modelde değil.
\cl{Tanım} \term{API}, eylemin teknik arabirimidir; \term{Authorization}, o
eylemi kimin yapabileceğinin sınırıdır.
\cl{Model} Model resepsiyonist; API kapı, Authorization ise anahtar. Model
müşteriye kendi anahtarını dağıtamaz.
\cl{Ne zaman} Tool calling'in arkasında; her dış eylemde ve veri erişiminde.
\cl{Tuzak} Yetkiyi ``kontrol et'' diye prompt'ta modele bırakmak; araçta/kodda uygulamamak ---
prompt injection bunu aşar.
\cl{İlişkili} \term{Tool Calling} · \term{MCP} · \term{Guardrails} \,{\scriptsize\color{softgray}(İlke \#6, \#15)}
